\chapter{Introduzione}
\label{chap:introduction}

La sicurezza alimentare e la vigilanza sui prodotti regolamentati costituiscono un ambito critico di interesse pubblico e industriale, in cui la tutela della salute dei consumatori si intreccia strettamente con le discipline del \textit{Project Management} e del \textit{Supply Chain Risk Management}. Lungo filiere agroalimentari, farmaceutiche e manifatturiere sempre più interconnesse su scala transnazionale, un'anomalia microbiologica, un'adulterazione chimica o un difetto di etichettatura possono propagarsi rapidamente dai siti di produzione e dai varchi doganali fino alla distribuzione al dettaglio, generando impatti sanitari, costi logistici di ritiro, sanzioni regolatorie e danni reputazionali durevoli.

In questo scenario, le autorità di regolazione producono e pubblicano archivi pubblici di straordinaria ricchezza informativa, ma strutturalmente eterogenei per granularità, modello di rilevazione e finalità di intervento. Dalla prospettiva del governo del rischio, l'osservazione isolata di un singolo registro fornisce una visione parziale della vulnerabilità delle filiere. Al contrario, la lettura congiunta di più fonti istituzionali consente di presidiare l'intero ciclo di vita del rischio regolatorio: dalla sorveglianza chimica preventiva sui campioni di laboratorio, al blocco delle partite non conformi alle frontiere, fino alla gestione post-vendita dei richiami sul mercato interno.

Le metodologie di \textit{Business Intelligence} (BI) e \textit{Visual Analytics} offrono il supporto tecnico e metodologico necessario per trasformare questi volumi di dati grezzi e frammentati in strumenti decisionali interattivi, governati da pipeline riproducibili di \textit{Extract, Transform, Load} (ETL), modellazione dimensionale e progettazione visiva orientata alla chiarezza quantitativa \citep{sharda2020,few2013}. Su queste premesse nasce il progetto \textbf{DashMe}, sviluppato nell'ambito del corso di \textit{Project Management \& Business Intelligence} presso l'Università Politecnica delle Marche (UnivPM), che integra tre moduli analitici distinti realizzati con le tre principali piattaforme di BI contemporanee, seguendo la sequenza logica e architettonica dal livello associativo esplorativo fino alla modellazione relazionale avanzata:

\begin{enumerate}
    \item \textbf{Qlik Sense -- Sorveglianza e Allerta Rapida Europea (RASFF)}: il primo pilastro analizza le notifiche del sistema europeo \textit{Rapid Alert System for Food and Feed} \citep{rasff_portal}. Sfruttando il motore associativo in-memory di Qlik Sense, il modulo esamina \textbf{11.787 notifiche validate} (2012--2025) provenienti da 30 Paesi notificanti verso 141 Paesi di origine, concentrandosi su urgenza operativa (43,73\% di allarmi immediati), gravità clinica (66,8\% di rischio grave), respingimenti alla frontiera (10,6\%, pari a 1.255 casi) e prevalenza biologica (\emph{Salmonella spp.} prima causa con 2.893 segnalazioni).

    \item \textbf{Tableau -- Richiami Regolatori Statunitensi (FDA)}: il secondo pilastro adotta l'ambiente di \textit{Visual Analytics} di Tableau per esaminare i rapporti di \textit{enforcement} della \textit{U.S. Food and Drug Administration} \citep{fda_recalls_portal}. A partire da \textbf{104.537 prodotti richiamati} associati a \textbf{37.848 eventi univoci} nell'arco 2012--2026 e arricchiti mediante 14 campi calcolati in lingua italiana, il modulo indaga la geografia industriale statunitense, i profili di severità clinica nelle tre classi regolatorie FDA (Classe~I, Classe~II e Classe~III) \citep{fda_recall_classes}, l'epidemiologia di cause e allergeni nonché la governance operativa dei 18.555 richiami tuttora \emph{In corso}, della catena del freddo e del confezionamento (ove il packaging in vetro registra un'incidenza del 74,4\% di adulterazioni farmacologiche).

    \item \textbf{Microsoft Power BI -- Sorveglianza Chimica Preventiva (EFSA / WHO GEMS)}: il terzo pilastro affronta la prevenzione analitica a monte tramite i dati europei dell'archivio congiunto dell'Autorità Europea per la Sicurezza Alimentare e dell'Organizzazione Mondiale della Sanità \citep{who_gems_efsa}. Basato su uno \textit{Star Schema} composto da una tabella dei fatti e quattro tabelle dimensionali su \textbf{2.424.475 rilevazioni laboratoristiche} (1984--2017), il modulo impiega misure DAX rigorose per determinare il tasso di positività effettivo al netto dei 3.330 valori nulli (\textbf{39,63\%} globale; 43,44\% sui 32 Stati nazionali nella prima dashboard; 36,96\% sugli analiti specifici nella seconda), evidenziando criticità di bioaccumulo (65.604 positivi al cadmio nella carne e nelle frattaglie) e identificando i driver statistici della contaminazione mediante algoritmi di intelligenza artificiale visiva (\emph{Key Influencers}, fino a 2,53x).
\end{enumerate}

La Tabella~\ref{tab:intro_modules_overview} sintetizza le caratteristiche strutturali, i volumi informativi e l'articolazione in viste interattive dei tre moduli che compongono l'ecosistema DashMe.

\begin{table}[htbp]
  \centering
  \caption{Quadro sinottico dei tre moduli analitici del progetto DashMe.}
  \label{tab:intro_modules_overview}
  \small
  \begin{tabularx}{\textwidth}{@{}>{\raggedright\arraybackslash}p{2.2cm}>{\raggedright\arraybackslash}X>{\raggedright\arraybackslash}Xp{1.7cm}p{1.9cm}@{}}
    \toprule
    \textbf{Piattaforma BI} & \textbf{Fonte Dati / Ente} & \textbf{Granularità e Volume} & \textbf{Arco Temp.} & \textbf{N. Dashboard} \\
    \midrule
    Qlik Sense &
    Portale RASFF / Commissione Europea (\texttt{RASFF\_\-window.csv}) &
    11.787 notifiche validate (30 Paesi notificanti, 141 Paesi di origine) &
    2012--2025 &
    4 fogli \\
    \addlinespace[0.4em]
    Tableau &
    FDA Enforcement Reports (\texttt{recalls.xlsx}) &
    104.537 prodotti in 37.848 eventi (14 campi calcolati) &
    2012--2026 &
    3 dashboard \\
    \addlinespace[0.4em]
    Power BI &
    EFSA / WHO GEMS (\texttt{Contamination\_\-european\_\-food.csv}) &
    2.424.475 rilevazioni (33 entità geografiche, 50 contaminanti, 348 alimenti; Star Schema 1 Fact + 4 Dim) &
    1984--2017 &
    3 dashboard \\
    \bottomrule
  \end{tabularx}
\end{table}

Alla luce di tale architettura multi-piattaforma, il presente elaborato si propone di rispondere a quattro domande di ricerca (\textit{Research Questions}, RQ):

\begin{enumerate}
    \item \textbf{RQ1 (Qlik Sense -- Sorveglianza e allerta rapida UE RASFF)}: quali trend storici, flussi geopolitici di confine, matrici alimentare-pericolo e dinamiche di tracciabilità emergono dalle notifiche RASFF?
    \item \textbf{RQ2 (Tableau -- Richiami regolatori USA FDA)}: come si articolano geografia industriale, profili di severità clinica (Classi I--III), epidemiologia delle cause/allergeni e governance operativa (catena del freddo, packaging, portata distributiva) nei richiami FDA?
    \item \textbf{RQ3 (Power BI -- Sorveglianza chimica europea EFSA / WHO GEMS)}: quali relazioni intercorrono tra volume di campionamento, tassi di positività al netto dei valori nulli, fasce di concentrazione per filiera, screening multi-analita e fattori di influenza AI?
    \item \textbf{RQ4 (Valutazione comparativa multi-piattaforma)}: in che misura Qlik Sense (motore associativo), Tableau (visual analytics esplorativa) e Power BI (modellazione stellare VertiPaq, DAX e AI visiva) offrono vantaggi e compromessi complementari?
\end{enumerate}

\section*{Struttura dell'Elaborato}
Il resto della trattazione è articolato in sei capitoli successivi:
\begin{itemize}
    \item il \textbf{Capitolo~\ref{chap:methodology}} illustra l'architettura metodologica complessiva del progetto DashMe, descrivendo le basi dati sorgente, le pipeline di \emph{data cleansing} e \emph{feature engineering} e i criteri di progettazione visiva adottati sulle tre piattaforme;
    \item il \textbf{Capitolo~\ref{chap:qlik}} presenta i risultati dell'analisi in \emph{Qlik Sense} sul sistema europeo di allerta rapida RASFF lungo quattro fogli dedicati a sintesi esecutiva, geopolitica dei controlli, matrici alimentari e tracciabilità di filiera;
    \item il \textbf{Capitolo~\ref{chap:tableau}} approfondisce le evidenze prodotte in \emph{Tableau} sui richiami regolatori statunitensi FDA, coprendo geografia industriale e responsabilità aziendale, epidemiologia delle cause con focus sugli allergeni e governance operativa su logistica e packaging;
    \item il \textbf{Capitolo~\ref{chap:powerbi}} esamina il modulo \emph{Microsoft Power BI} per la sorveglianza chimica europea EFSA / WHO GEMS, analizzando il paradigma volume-positività, il profilo di rischio per filiera e fasce di concentrazione, le dinamiche multi-analita e i fattori determinanti individuati tramite AI visiva;
    \item il \textbf{Capitolo~\ref{chap:discussion}} propone una discussione trasversale delle implicazioni manageriali per il governo delle filiere, confronta criticamente i tre strumenti di BI rispetto alle capacità tecniche osservate e discute i limiti metodologici dei dataset analizzati;
    \item il \textbf{Capitolo~\ref{chap:conclusions}} sintetizza le risposte alle quattro domande di ricerca e delinea le possibili estensioni future dell'ecosistema DashMe.
\end{itemize}
